\chapter{Introduction}\label{ch:intro}

\section{Motivation}\label{sec:motivation}

Trie-based index structures are the backbone of modern in-memory databases.
Since the Adaptive Radix Tree (ART) by Leis et al.~\cite{leis2013art} in 2013,
research has diversified in several directions: trie-height optimization
(HOT~\cite{binna2018hot}), hybrid B\textsuperscript{+} structures
(Masstree~\cite{mao2012masstree}, B\textsuperscript{2}-tree~\cite{schmeisser2022b2tree}),
succinct representations (CoCo-Trie~\cite{boffa2024coco}, SuRF~\cite{zhang2018surf}),
and self-tuning approaches (START~\cite{fent2020start}).

In parallel, a dedicated body of research on cache awareness, prefetching, and
allocation strategies has emerged. Heuristics such as Cache-Sensitive Search
Trees, cache-oblivious layouts, and hierarchical prefetching each address an
individual axis of the cache hierarchy --- yet they have never been
systematically combined within a single, permuted measurement architecture.

This thesis closes that gap with a three-layer cache-engine architecture
(CacheEngine $\to$ ExecutionEngine $\to$ SearchEngine) whose building blocks are
decomposed into fourteen orthogonal permutation axes, spanning a configuration
space of more than $10^{11}$ permutations. As a probe for state-of-the-art
membership we contribute PRT-ART (a Probst Redirect Tree / ART hybrid) with eight
building-block layers of its own.

The framework rests on an anatomical metaphor: a search algorithm is an
\emph{animal} (organism) whose \emph{organs} are the axes --- the sub-tasks every
algorithm must perform (node layout, traversal, allocation, prefetch, and so on),
each in a different expression. A concrete algorithm is then one \emph{composition}
of organ expressions, a single point in the Cartesian product of all axes; permuting
the axes is a genetic experiment that breeds new composite algorithms with previously
unknown cache-line behaviour. The research goal is to find this shared anatomy and the
fastest ``animal'' for a given data-type-favourable load pattern.

\section{Problem Statement}\label{sec:problem}

Existing cache-aware index designs optimize one or a few aspects in isolation
and report results that are hard to compare across publications, because each
work uses its own data structures, workloads, and measurement methodology. There
is no common framework that (i) decomposes such designs into interchangeable
axes, (ii) reconstructs published algorithms as explicit configurations, and
(iii) measures both individual axes and whole algorithms on a single, fair basis.

\section{Research Questions}\label{sec:rqs}

\begin{description}
  \item[RQ1 (Composition)] Which algorithm axes of a cache-aware trie index can
        be canonically systematized in a permutation matrix, such that
        state-of-the-art implementations can be reconstructed as profile
        configurations?
  \item[RQ2 (Measurement methodology)] How can a fair, reproducible YCSB workload
        run be performed per permutation, allowing profile-specific workload
        selection (e.g.\ a range-scan profile $\to$ YCSB-E) while minimizing the
        measurement overhead via an opt-in experiment mode?
  \item[RQ3 (Probe comparison)] How does PRT-ART perform as a probe against the
        eight Tier-1 SOTA designs (ART, HOT, Masstree, CoCo-Trie, START,
        B\textsuperscript{2}-tree, Wormhole, SuRF) and the Tier-2/3 SOTA under the
        three mandatory measurement series A (probe vs.\ SOTA), B (SOTA-only
        baseline), and C (merge points)?
\end{description}

\section{Objectives and Contributions}\label{sec:contributions}

The central contributions of this thesis are:

\begin{itemize}
  \item A three-layer architecture (CacheEngine $\to$ ExecutionEngine $\to$
        SearchEngine) with ABI-stable C++23 module interfaces and variadically
        templated hybrid-API specializations (1~parameter $\Rightarrow$
        \texttt{std::vector} API; 2~parameters $\Rightarrow$ \texttt{std::map}
        API; $N>2 \Rightarrow$ \texttt{std::map<K, std::tuple<V$_1$\ldots V$_N$>>}).
  \item A fourteen-axis permutation matrix with \texttt{std::variant} building
        blocks and compile-time fallback between the probe stack and the
        state-of-the-art stack.
  \item Thirty SOTA algorithm profiles as persisted XML configurations (8~Tier-1
        + 22~Tier-2/3) that the \texttt{ExperimentDriver} transforms automatically
        into pre-compiled C++23 modules.
  \item PRT-ART as a probe with eight building-block layers of its own (4+2 pool
        allocator, optimistic lock coupling with reserved blocks, multi-level
        layout, path-oriented prefetch, density tracker, hypothesis metrics
        H1/H2/H3, inline/external/chain-ref value handles, signaling-bit
        serialization).
  \item A measurement driver with defined-/full-mode distinction and
        profile-aware workload routing.
\end{itemize}

\section{Structure of the Thesis}\label{sec:structure}

Chapter~\ref{ch:fundamentals} introduces the necessary background. Chapter~\ref{ch:sota}
summarizes the state of the art across six search clusters (A--F) and an
allocator cluster. Chapter~\ref{ch:concept} presents the axis library framework,
the layered architecture, the ABI interface, and the fourteen-axis permutation
matrix. Chapter~\ref{ch:impl} documents the implementation across the three
repositories. Chapter~\ref{ch:methodology} describes the measurement methodology
with the three mandatory series and profile-aware workload routing.
Chapter~\ref{ch:results} presents and discusses the results, including the
hypotheses H1--H3. Chapter~\ref{ch:conclusion} concludes with a summary and an
outlook.
